We consider the Wilson lattice gauge measure for $\SU$ on a periodic three-dimensional torus of fixed size, with bare coupling $g^2=g_3^2a$ proportional to the lattice spacing $a$, as dictated by the superrenormalizability of three-dimensional Yang--Mills theory. Observables are Wilson loops of the lattice configuration evolved by the Wilson flow, the lattice gradient flow of the Wilson action, at positive flow time. We prove that, at every fixed value of the coupling, the joint laws of finitely many flowed Wilson loops converge as the lattice spacing tends to zero along the sequences of lattice sides $nM^k$, for fixed odd $M\ge3$ and fixed positive $n$ divisible by four. For each such refinement class, the limits are realized by a unique probability measure on the space of gauge orbits of Yang--Mills heat-flow trajectories at positive times, a version of the state space introduced by Cao and Chatterjee. Under this measure the energy at positive times is tight, and the measure is invariant under the lattice symmetries and the center symmetries. We also show that, at short distances, the rescaled deficits $(1-W)/(g_3^2\varepsilon)$ of the normalized fundamental Wilson loops $W$ of loops $x+\varepsilon\ell$, flowed for time $\sigma\varepsilon^2$ with fixed $\sigma>0$, converge to explicit quadratic functions of the heat-smoothed Gaussian free field. As the coupling tends to zero, the measure concentrates on flat connections, with Gaussian fluctuations of explicit covariance near regular commuting holonomies and fluctuations governed by the constant-mode quartic form near the central holonomies. The proof is a multiscale renormalization group analysis in the tradition of Bałaban, carried out on a small-field event whose complement we show to have small probability, uniformly in the lattice spacing, using reflection positivity and chessboard estimates. It is combined with stochastic estimates for the lattice heat flow and a deterministic analysis of the Wilson flow for rough initial data.
